\Needspace{8\baselineskip}\subsection{\texorpdfstring{\nolinkurl{rnj_wzzt/cli.py}}{rnj\_wzzt/cli.py}}\label{module:cli}
普通入口配置与高级参数解析。所有行号对应随附源码快照。
\begin{longtable}{P{0.73\textwidth}P{0.18\textwidth}}\toprule 符号（类或函数） & 原始行号\\\midrule\endhead
\nolinkurl{_add_scenario_options} & 15--20 \\
\nolinkurl{_scenario_options} & 23--27 \\
\nolinkurl{run} & 30--69 \\
\nolinkurl{main} & 72--86 \\
\nolinkurl{advanced_main} & 89--135 \\
\bottomrule\end{longtable}
\Needspace{8\baselineskip}\subsection{\texorpdfstring{\nolinkurl{rnj_wzzt/data/paper_style_case_bank.py}}{rnj\_wzzt/data/paper\_style\_case\_bank.py}}\label{module:data.paper_style_case_bank}
四算例注册表、三种网络及资源配置。所有行号对应随附源码快照。
\begin{longtable}{P{0.73\textwidth}P{0.18\textwidth}}\toprule 符号（类或函数） & 原始行号\\\midrule\endhead
\nolinkurl{_z} & 13--14 \\
\nolinkurl{_make_case} & 17--103 \\
\nolinkurl{build_soumalas_style_lv_11_case} & 106--147 \\
\nolinkurl{build_flynn_balanced_lv_16_case} & 150--182 \\
\nolinkurl{build_pengwah_mixed_lv_18_case} & 185--211 \\
\nolinkurl{build_case_bank_resource_assignment} & 221--260 \\
\bottomrule\end{longtable}
\Needspace{8\baselineskip}\subsection{\texorpdfstring{\nolinkurl{rnj_wzzt/data/paper_style_terminal_lv.py}}{rnj\_wzzt/data/paper\_style\_terminal\_lv.py}}\label{module:data.paper_style_terminal_lv}
paper15 网络及异质终端资源配置。所有行号对应随附源码快照。
\begin{longtable}{P{0.73\textwidth}P{0.18\textwidth}}\toprule 符号（类或函数） & 原始行号\\\midrule\endhead
\nolinkurl{build_paper_style_terminal_lv_case} & 17--156 \\
\nolinkurl{build_paper_style_resource_assignment} & 159--203 \\
\bottomrule\end{longtable}
\Needspace{8\baselineskip}\subsection{\texorpdfstring{\nolinkurl{rnj_wzzt/estimation/constrained_least_squares.py}}{rnj\_wzzt/estimation/constrained\_least\_squares.py}}\label{module:estimation.constrained_least_squares}
对称非负有序 R/X 凸 QP 数值求解。所有行号对应随附源码快照。
\begin{longtable}{P{0.73\textwidth}P{0.18\textwidth}}\toprule 符号（类或函数） & 原始行号\\\midrule\endhead
\nolinkurl{make_feasible} & 10--20 \\
\nolinkurl{solve_symmetric_least_squares} & 23--215 \\
\bottomrule\end{longtable}
\Needspace{8\baselineskip}\subsection{\texorpdfstring{\nolinkurl{rnj_wzzt/estimation/laminar_l1_milp.py}}{rnj\_wzzt/estimation/laminar\_l1\_milp.py}}\label{module:estimation.laminar_l1_milp}
固定支持 LP、一步 MILP、路径与验证选择。所有行号对应随附源码快照。
\begin{longtable}{P{0.73\textwidth}P{0.18\textwidth}}\toprule 符号（类或函数） & 原始行号\\\midrule\endhead
\nolinkurl{SolverDiagnostics} & 40--56 \\
\nolinkurl{SolverDiagnostics.to_dict} & 55--56 \\
\nolinkurl{FixedSupportSolution} & 60--68 \\
\nolinkurl{ExtensionSolution} & 72--80 \\
\nolinkurl{LaminarL1PathPoint} & 84--99 \\
\nolinkurl{LaminarL1Result} & 103--139 \\
\nolinkurl{LaminarL1Result.summary} & 124--139 \\
\nolinkurl{_PreparedScenarios} & 143--160 \\
\nolinkurl{_PreparedScenarios.n} & 151--152 \\
\nolinkurl{_PreparedScenarios.scenario_count} & 155--156 \\
\nolinkurl{_PreparedScenarios.observation_count} & 159--160 \\
\nolinkurl{_require_positive_finite} & 163--167 \\
\nolinkurl{_require_nonnegative_finite} & 170--174 \\
\nolinkurl{_validate_solver_parameters} & 177--196 \\
\nolinkurl{solver_diagnostics_prove_optimality} & 199--229 \\
\nolinkurl{_dual_bound_certifies_no_gain} & 232--244 \\
\nolinkurl{_objectives_agree} & 247--252 \\
\nolinkurl{_VariableBuilder} & 255--289 \\
\nolinkurl{_VariableBuilder.__init__} & 258--263 \\
\nolinkurl{_VariableBuilder.add} & 265--285 \\
\nolinkurl{_VariableBuilder.size} & 288--289 \\
\nolinkurl{_ConstraintBuilder} & 292--332 \\
\nolinkurl{_ConstraintBuilder.__init__} & 295--300 \\
\nolinkurl{_ConstraintBuilder.add} & 302--316 \\
\nolinkurl{_ConstraintBuilder.build} & 318--328 \\
\nolinkurl{_ConstraintBuilder.size} & 331--332 \\
\nolinkurl{_matrix_from_scenario} & 335--345 \\
\nolinkurl{_prepare_scenarios} & 348--388 \\
\nolinkurl{normalize_supports} & 391--410 \\
\nolinkurl{normalize_support_pool} & 413--434 \\
\nolinkurl{is_laminar_family} & 437--446 \\
\nolinkurl{is_admissible_extension} & 449--460 \\
\nolinkurl{_round_and_validate_support} & 463--477 \\
\nolinkurl{build_matrices_from_atoms} & 480--522 \\
\nolinkurl{_fixed_atom_features} & 525--543 \\
\nolinkurl{_target_and_scenario_output} & 546--553 \\
\nolinkurl{_diagnostics} & 556--578 \\
\nolinkurl{_run_milp} & 581--618 \\
\nolinkurl{_add_absolute_residual_constraints} & 621--673 \\
\nolinkurl{_extract_intercepts} & 676--679 \\
\nolinkurl{_solve_fixed_prepared} & 682--740 \\
\nolinkurl{solve_fixed_support_l1} & 743--770 \\
\nolinkurl{_upper_pairs} & 773--775 \\
\nolinkurl{_solve_extension_prepared} & 778--979 \\
\nolinkurl{solve_best_laminar_extension_l1} & 982--1023 \\
\nolinkurl{estimate_atom_upper_bounds} & 1026--1054 \\
\nolinkurl{evaluate_l1_matrices} & 1057--1113 \\
\nolinkurl{_labels_for_supports} & 1116--1119 \\
\nolinkurl{_make_path_point} & 1122--1162 \\
\nolinkurl{_touches_bound} & 1165--1166 \\
\nolinkurl{_choose_path_point} & 1169--1179 \\
\nolinkurl{fit_laminar_l1_sensitivity} & 1182--1460 \\
\bottomrule\end{longtable}
\Needspace{8\baselineskip}\subsection{\texorpdfstring{\nolinkurl{rnj_wzzt/estimation/matrix_constraints.py}}{rnj\_wzzt/estimation/matrix\_constraints.py}}\label{module:estimation.matrix_constraints}
矩阵可行性修正和树度量诊断。所有行号对应随附源码快照。
\begin{longtable}{P{0.73\textwidth}P{0.18\textwidth}}\toprule 符号（类或函数） & 原始行号\\\midrule\endhead
\nolinkurl{SensitivityMatrixDiagnostics} & 11--26 \\
\nolinkurl{SensitivityMatrixDiagnostics.to_dict} & 23--26 \\
\nolinkurl{_matrix_scale} & 29--34 \\
\nolinkurl{_project_diagonal_order} & 37--62 \\
\nolinkurl{_make_ordered_feasible} & 65--76 \\
\nolinkurl{project_ordered_sensitivity_matrix} & 79--117 \\
\nolinkurl{project_tree_covariance_matrix} & 121--183 \\
\nolinkurl{sensitivity_matrix_diagnostics} & 186--219 \\
\nolinkurl{four_point_violation_summary} & 222--248 \\
\bottomrule\end{longtable}
\Needspace{8\baselineskip}\subsection{\texorpdfstring{\nolinkurl{rnj_wzzt/estimation/multiscenario.py}}{rnj\_wzzt/estimation/multiscenario.py}}\label{module:estimation.multiscenario}
标签对齐、消截距、回归接口与诊断。所有行号对应随附源码快照。
\begin{longtable}{P{0.73\textwidth}P{0.18\textwidth}}\toprule 符号（类或函数） & 原始行号\\\midrule\endhead
\nolinkurl{preprocess_scenarios} & 16--30 \\
\nolinkurl{_aligned_arrays} & 33--59 \\
\nolinkurl{_symmetric_blocks} & 62--67 \\
\nolinkurl{_fixed_margins} & 70--78 \\
\nolinkurl{_fit_tree_covariance_compatibility} & 81--116 \\
\nolinkurl{fit_projected_sensitivity} & 119--204 \\
\bottomrule\end{longtable}
\Needspace{8\baselineskip}\subsection{\texorpdfstring{\nolinkurl{rnj_wzzt/estimation/preprocessing.py}}{rnj\_wzzt/estimation/preprocessing.py}}\label{module:estimation.preprocessing}
平方电压降、去均值及其他时序变换。所有行号对应随附源码快照。
\begin{longtable}{P{0.73\textwidth}P{0.18\textwidth}}\toprule 符号（类或函数） & 原始行号\\\midrule\endhead
\nolinkurl{squared_voltage_drop_from_observed_root} & 11--29 \\
\nolinkurl{daily_demean} & 32--46 \\
\nolinkurl{rolling_highpass} & 49--56 \\
\nolinkurl{apply_preprocessing_recipe} & 59--80 \\
\bottomrule\end{longtable}
\Needspace{8\baselineskip}\subsection{\texorpdfstring{\nolinkurl{rnj_wzzt/graph/bootstrap.py}}{rnj\_wzzt/graph/bootstrap.py}}\label{module:graph.bootstrap}
循环分块重采样、边界块及不相交筛选。所有行号对应随附源码快照。
\begin{longtable}{P{0.73\textwidth}P{0.18\textwidth}}\toprule 符号（类或函数） & 原始行号\\\midrule\endhead
\nolinkurl{_boundary_cherries} & 11--18 \\
\nolinkurl{_moving_block_bootstrap_copy} & 21--41 \\
\nolinkurl{_select_disjoint} & 44--60 \\
\bottomrule\end{longtable}
\Needspace{8\baselineskip}\subsection{\texorpdfstring{\nolinkurl{rnj_wzzt/graph/rooted_hierarchy.py}}{rnj\_wzzt/graph/rooted\_hierarchy.py}}\label{module:graph.rooted_hierarchy}
clade 转换、伪终端反嵌入、局部重辨识工具。所有行号对应随附源码快照。
\begin{longtable}{P{0.73\textwidth}P{0.18\textwidth}}\toprule 符号（类或函数） & 原始行号\\\midrule\endhead
\nolinkurl{PseudoCluster} & 19--26 \\
\nolinkurl{rooted_tree_from_clades} & 29--87 \\
\nolinkurl{rooted_clades} & 90--117 \\
\nolinkurl{select_peripheral_clusters} & 120--185 \\
\nolinkurl{_boundary_sensitivity_row} & 188--208 \\
\nolinkurl{aggregate_rooted_scenarios} & 211--308 \\
\nolinkurl{build_local_cluster_scenarios} & 311--346 \\
\nolinkurl{reidentify_local_clades_from_shared_paths} & 349--392 \\
\nolinkurl{expand_pseudo_clades} & 395--410 \\
\nolinkurl{expand_pseudo_clades_with_local_reidentification} & 413--430 \\
\nolinkurl{expand_pseudo_sibling_pairs} & 433--456 \\
\bottomrule\end{longtable}
\Needspace{8\baselineskip}\subsection{\texorpdfstring{\nolinkurl{rnj_wzzt/graph/rooted_neighbor_joining.py}}{rnj\_wzzt/graph/rooted\_neighbor\_joining.py}}\label{module:graph.rooted_neighbor_joining}
按共享路径递归合并的一般树 RNJ。所有行号对应随附源码快照。
\begin{longtable}{P{0.73\textwidth}P{0.18\textwidth}}\toprule 符号（类或函数） & 原始行号\\\midrule\endhead
\nolinkurl{RootedTreeResult} & 21--27 \\
\nolinkurl{shared_paths_from_distances} & 30--46 \\
\nolinkurl{_contract_zero_internal_edges} & 49--85 \\
\nolinkurl{rooted_neighbor_joining} & 88--178 \\
\bottomrule\end{longtable}
\Needspace{8\baselineskip}\subsection{\texorpdfstring{\nolinkurl{rnj_wzzt/graph/sensitivity_geometry.py}}{rnj\_wzzt/graph/sensitivity\_geometry.py}}\label{module:graph.sensitivity_geometry}
R/X 距离归一化及 RNJ 共享路径输入。所有行号对应随附源码快照。
\begin{longtable}{P{0.73\textwidth}P{0.18\textwidth}}\toprule 符号（类或函数） & 原始行号\\\midrule\endhead
\nolinkurl{SensitivityGeometry} & 26--35 \\
\nolinkurl{_validated_matrix} & 38--44 \\
\nolinkurl{_positive_mean} & 47--49 \\
\nolinkurl{_mode_coefficients} & 52--60 \\
\nolinkurl{distance_candidates} & 63--81 \\
\nolinkurl{sensitivity_geometry} & 84--123 \\
\bottomrule\end{longtable}
\Needspace{8\baselineskip}\subsection{\texorpdfstring{\nolinkurl{rnj_wzzt/models/ac_powerflow.py}}{rnj\_wzzt/models/ac\_powerflow.py}}\label{module:models.ac_powerflow}
径向 AC 后向前向扫掠潮流。所有行号对应随附源码快照。
\begin{longtable}{P{0.73\textwidth}P{0.18\textwidth}}\toprule 符号（类或函数） & 原始行号\\\midrule\endhead
\nolinkurl{ACPowerFlowSnapshot} & 22--32 \\
\nolinkurl{_oriented_tree} & 35--57 \\
\nolinkurl{solve_ac_power_flow_snapshot} & 60--95 \\
\nolinkurl{_solve_ac_power_flow_snapshot_oriented} & 98--168 \\
\nolinkurl{solve_ac_power_flow_timeseries} & 171--292 \\
\bottomrule\end{longtable}
\Needspace{8\baselineskip}\subsection{\texorpdfstring{\nolinkurl{rnj_wzzt/models/lin_distflow.py}}{rnj\_wzzt/models/lin\_distflow.py}}\label{module:models.lin_distflow}
标幺换算、理想共享路径 R/X 与距离。所有行号对应随附源码快照。
\begin{longtable}{P{0.73\textwidth}P{0.18\textwidth}}\toprule 符号（类或函数） & 原始行号\\\midrule\endhead
\nolinkurl{ohm_to_pu} & 11--15 \\
\nolinkurl{build_reduced_sensitivity_matrices} & 18--60 \\
\nolinkurl{impedance_distance_from_reduced_R} & 63--67 \\
\nolinkurl{impedance_distance_from_reduced_X} & 70--73 \\
\bottomrule\end{longtable}
\Needspace{8\baselineskip}\subsection{\texorpdfstring{\nolinkurl{rnj_wzzt/models/network.py}}{rnj\_wzzt/models/network.py}}\label{module:models.network}
网络对象、图转换与校验入口。所有行号对应随附源码快照。
\begin{longtable}{P{0.73\textwidth}P{0.18\textwidth}}\toprule 符号（类或函数） & 原始行号\\\midrule\endhead
\nolinkurl{TerminalizedNetwork} & 12--101 \\
\nolinkurl{TerminalizedNetwork.observed_buses} & 24--27 \\
\nolinkurl{TerminalizedNetwork.injection_buses} & 29--32 \\
\nolinkurl{TerminalizedNetwork.hidden_buses} & 34--37 \\
\nolinkurl{TerminalizedNetwork.load_buses} & 39--42 \\
\nolinkurl{TerminalizedNetwork.closed_edges} & 44--48 \\
\nolinkurl{TerminalizedNetwork.candidate_edges} & 50--58 \\
\nolinkurl{TerminalizedNetwork.to_networkx_graph} & 60--74 \\
\nolinkurl{TerminalizedNetwork.orient_from_root} & 76--85 \\
\nolinkurl{TerminalizedNetwork.get_leaf_buses} & 87--94 \\
\nolinkurl{TerminalizedNetwork.check_terminal_load_only} & 96--101 \\
\bottomrule\end{longtable}
\Needspace{8\baselineskip}\subsection{\texorpdfstring{\nolinkurl{rnj_wzzt/pipeline.py}}{rnj\_wzzt/pipeline.py}}\label{module:pipeline}
数据准备、RNJ、收缩、MILP、输出的流程编排。所有行号对应随附源码快照。
\begin{longtable}{P{0.73\textwidth}P{0.18\textwidth}}\toprule 符号（类或函数） & 原始行号\\\midrule\endhead
\nolinkurl{_CaseData} & 53--62 \\
\nolinkurl{_RnjSelection} & 66--74 \\
\nolinkurl{_validate_run_options} & 77--120 \\
\nolinkurl{_scenario_manifest_rows} & 123--143 \\
\nolinkurl{_rx75_rnj_clades} & 146--166 \\
\nolinkurl{_select_boundary_blocks} & 169--211 \\
\nolinkurl{_align_scenario_names} & 214--222 \\
\nolinkurl{_prepare_case} & 225--289 \\
\nolinkurl{_select_case_rnj} & 292--332 \\
\nolinkurl{_contracted_milp_inputs} & 335--397 \\
\nolinkurl{_leaf_singletons} & 400--405 \\
\nolinkurl{_map_clade_to_reduced_support} & 408--431 \\
\nolinkurl{_rnj_reduced_candidate_pool} & 434--463 \\
\nolinkurl{_expand_candidate_pool} & 466--479 \\
\nolinkurl{_expand_pseudo_result_clades} & 482--492 \\
\nolinkurl{_fit_milp_variants} & 495--593 \\
\nolinkurl{run} & 596--755 \\
\nolinkurl{main} & 758--762 \\
\bottomrule\end{longtable}
\Needspace{8\baselineskip}\subsection{\texorpdfstring{\nolinkurl{rnj_wzzt/reporting.py}}{rnj\_wzzt/reporting.py}}\label{module:reporting}
真值评价、候选审计与求解器记录。所有行号对应随附源码快照。
\begin{longtable}{P{0.73\textwidth}P{0.18\textwidth}}\toprule 符号（类或函数） & 原始行号\\\midrule\endhead
\nolinkurl{_truth_nontrivial_clades} & 9--19 \\
\nolinkurl{_family_score} & 22--34 \\
\nolinkurl{_serialize} & 37--38 \\
\nolinkurl{_serialize_family} & 41--45 \\
\nolinkurl{_rnj_summary_row} & 48--94 \\
\nolinkurl{_rnj_candidate_rows} & 97--130 \\
\nolinkurl{_milp_result_row} & 133--227 \\
\bottomrule\end{longtable}
\Needspace{8\baselineskip}\subsection{\texorpdfstring{\nolinkurl{rnj_wzzt/scenario/profiles.py}}{rnj\_wzzt/scenario/profiles.py}}\label{module:scenario.profiles}
负荷、光伏、风电、小型机组及无功曲线。所有行号对应随附源码快照。
\begin{longtable}{P{0.73\textwidth}P{0.18\textwidth}}\toprule 符号（类或函数） & 原始行号\\\midrule\endhead
\nolinkurl{_daily_demand_multiplier} & 9--24 \\
\nolinkurl{_pv_profile} & 27--43 \\
\nolinkurl{_wind_profile} & 46--59 \\
\nolinkurl{_small_generator_profile} & 62--80 \\
\nolinkurl{_curve_description} & 83--91 \\
\nolinkurl{_customer_type} & 94--101 \\
\nolinkurl{_resource_components} & 104--125 \\
\nolinkurl{_demand_power_factor} & 128--135 \\
\nolinkurl{_q_from_pf} & 138--142 \\
\nolinkurl{_der_q_injection} & 145--153 \\
\nolinkurl{_scenario_shape} & 156--205 \\
\nolinkurl{_build_terminal_profiles} & 208--324 \\
\bottomrule\end{longtable}
\Needspace{8\baselineskip}\subsection{\texorpdfstring{\nolinkurl{rnj_wzzt/scenario/settings.py}}{rnj\_wzzt/scenario/settings.py}}\label{module:scenario.settings}
物理工况、根观测与量测噪声独立配置。所有行号对应随附源码快照。
\begin{longtable}{P{0.73\textwidth}P{0.18\textwidth}}\toprule 符号（类或函数） & 原始行号\\\midrule\endhead
\nolinkurl{_finite_number} & 12--22 \\
\nolinkurl{resolve_scenario_settings} & 25--78 \\
\bottomrule\end{longtable}
\Needspace{8\baselineskip}\subsection{\texorpdfstring{\nolinkurl{rnj_wzzt/scenario/simulation.py}}{rnj\_wzzt/scenario/simulation.py}}\label{module:scenario.simulation}
AC 仿真、噪声、训练验证复制与子采样。所有行号对应随附源码快照。
\begin{longtable}{P{0.73\textwidth}P{0.18\textwidth}}\toprule 符号（类或函数） & 原始行号\\\midrule\endhead
\nolinkurl{_terminal_buses} & 32--33 \\
\nolinkurl{_resource_assignment} & 36--39 \\
\nolinkurl{_integer} & 42--45 \\
\nolinkurl{_root_voltage} & 48--73 \\
\nolinkurl{_dynamic_rms} & 76--78 \\
\nolinkurl{_diagnostics} & 81--118 \\
\nolinkurl{_simulate_pool} & 121--260 \\
\bottomrule\end{longtable}
\Needspace{8\baselineskip}\subsection{\texorpdfstring{\nolinkurl{rnj_wzzt/scenario/validate_scenario.py}}{rnj\_wzzt/scenario/validate\_scenario.py}}\label{module:scenario.validate_scenario}
径向、末端负荷、可观测性与守恒检查。所有行号对应随附源码快照。
\begin{longtable}{P{0.73\textwidth}P{0.18\textwidth}}\toprule 符号（类或函数） & 原始行号\\\midrule\endhead
\nolinkurl{validate_terminal_load_only_scenario} & 13--97 \\
\bottomrule\end{longtable}